\input{../tex-inputs/latex-leseansicht-vorspann}

\section[Arthur Schnitzler und Hugo von Hofmannsthal an Gustav Schwarzkopf, 3. 7. {[}1902{]}]{L04441 Arthur Schnitzler und Hugo von Hofmannsthal an Gustav Schwarzkopf, 3. 7. {[}1902{]}}
\nopagebreak\mylabel{L04441v}
\rehead{ }\normalsize\beginnumbering\briefempfaengerindex{Schwarzkopf, Gustav@\textsc{Schwarzkopf, Gustav}!zzzHofmannsthal, Hugo von@\emph{von Hugo von Hofmannsthal}!1902-07-033@{3. 7. {[}1902{]}}|(be}\briefempfaengerindex{Schwarzkopf, Gustav@\textsc{Schwarzkopf, Gustav}!zzzSchnitzler, Arthur@\emph{von Arthur Schnitzler}!1902-07-033@{3. 7. {[}1902{]}}|(be}
\toendnotes[C]{\smallbreak\pagebreak[2]}
\correspDesc{Versand  durch Arthur Schnitzler, Hofmannsthal am 3. 7. [1902] in Matrei am Brenner
\newline{}Übermittlung  in Deutsch-Matrei
\newline{}Erhalt  durch Gustav Schwarzkopf im Zeitraum [4. 7. 1902 – 8. 7. 1902?] in Tiefer Graben 23}\toendnotes[C]{\smallbreak}
\Standort{DLA, A:Schnitzler, HS.1985.1.1897.}
\physDesc{Bildpostkarte, 83 Zeichen
\newline{}Handschrift Arthur Schnitzler: Bleistift, deutsche Kurrent
\newline{}Handschrift Hugo von Hofmannsthal: Bleistift
\newline{}Versand: Stempel: »\nobreak{}\oindex{Matrei am Brenner@\textbf{Matrei am Brenner}|pwk}Deutsch-\textcolor{gray}{Matrei}\nobreak{}«.  }\pstart{}{\pb}Hrn Guſtav Schwarzkopf\pend{}\pstart{}Wien\oindex{Wien@\textbf{Wien}|pw}\pend{}\pstart{}I. Tiefer Graben 23\oindex{Wien@\textbf{Wien}!I., Innere Stadt@\textbf{I., Innere Stadt}!Tiefer Graben 23@\textbf{Tiefer Graben 23}|pw}.\pend{}{\bigskip}
\pstart
           \centering{}{\pb}\textcolor{gray}{\textbf{MATREI\oindex{Matrei am Brenner@\textbf{Matrei am Brenner}|pw}.}}\pend
           \vspace{1em}
\pstart
           \noindent{}{\pb}Herzlichen Gruß!\pend
           \pstart Ihr \spacefill\mbox{Arthur}\pend{}\selectlanguage{ngerman}\vspace{1em}\pstart \spacefill\mbox{{[}hs. Hofmannsthal:{]} Hugo} \pend{}
\pstart
           \strikeout{1}3 VII.\pend
           \selectlanguage{ngerman}\endnumbering\briefempfaengerindex{Schwarzkopf, Gustav@\textsc{Schwarzkopf, Gustav}!zzzHofmannsthal, Hugo von@\emph{von Hugo von Hofmannsthal}!1902-07-033@{3. 7. {[}1902{]}}|)be}\briefempfaengerindex{Schwarzkopf, Gustav@\textsc{Schwarzkopf, Gustav}!zzzSchnitzler, Arthur@\emph{von Arthur Schnitzler}!1902-07-033@{3. 7. {[}1902{]}}|)be}\mylabel{L04441h}
\begin{anhang}
\end{anhang}\newcommand{\dateiname}{L04441}\newcommand{\titel}{Arthur Schnitzler und Hugo von Hofmannsthal an Gustav Schwarzkopf, 3. 7. [1902]}\newcommand{\editorInnen}{Herausgegeben von Jahnke, SelmaMüller, Martin Anton}\input{../tex-inputs/latex-leseansicht-abspann}
